% \documentclass[a4, 10pt]{scrartcl}
\documentclass[a4paper, 10pt, twocolumn]{amsart}

\usepackage[
% fancytheorems,
% fancyproofs
nokoma
]{adam}

\usepackage[margin=0.75in]{geometry}

\title{Principles of Quantum Mechanics}
\author{Adam Kelly -- Mathematical Tripos Part II}
\date{\today. Email \texttt{ak2316@cam.ac.uk}}

\begin{document}

\maketitle

\section{Uncertainty Principle}

For a normalized vector $\psi$ in a finite-dimensional complex inner product space and Hermitian operators $A,B$, write $\langle T\rangle_\psi=\langle\psi,T\psi\rangle$, using an inner product linear in its second argument. Define
\[
A_0=A-\langle A\rangle_\psi I,\qquad B_0=B-\langle B\rangle_\psi I,
\qquad (\Delta A)_\psi=\|A_0\psi\|,\quad (\Delta B)_\psi=\|B_0\psi\|.
\]
With $z=\langle A_0\psi,B_0\psi\rangle$, Hermiticity gives
$\langle[A,B]\rangle_\psi=z-\overline z=2i\operatorname{Im}z$.
Cauchy--Schwarz therefore yields
\[
(\Delta A)_\psi(\Delta B)_\psi\geq\tfrac12|\langle[A,B]\rangle_\psi|.
\]
This proof includes zero-variance cases. Noncommuting operators alone do not imply a strictly positive lower bound in every state: the expectation of the commutator can vanish.

For unbounded operators, domain hypotheses are essential. The same argument applies to a normalized vector $\psi$ in a common domain where $A\psi,B\psi,AB\psi,BA\psi$ exist, and where symmetry justifies the displayed inner-product identities. In particular, for position and momentum satisfying $[x,p]\psi=i\hbar\psi$ on such a domain, it gives
\[
(\Delta x)_\psi(\Delta p)_\psi\geq\hbar/2.
\]
This is a conditional application, not an assertion that arbitrary unbounded operators share such a domain. Canonical position and momentum are not finite-dimensional matrices.

\end{document}
